\documentclass[12pt]{article}
\usepackage{graphicx}
\usepackage{amsmath}
\usepackage{geometry}
\geometry{a4paper, margin=1in}
\usepackage{titlesec}
\usepackage{minted}

% Format sections to match the reference document style
\titleformat{\section}
  {\normalfont\Large\bfseries}{\thesection.}{1em}{}
\titleformat{\subsection}
  {\normalfont\large\bfseries}{\thesubsection}{1em}{}

\begin{document}

\begin{center}
    {\Large \textbf{Quantum Simulation Lab}} \\
    \vspace{0.8cm}
    {\Large Experiment IV} \\
    \vspace{0.2cm}
    {\Large \textbf{Quantum Fourier Transform}} \\
\end{center}

\vspace{0.5cm}
\noindent Name: Hrishikesh M Jayakrishnan \hfill ID: SC26D054
\vspace{1em}
\hrule

\vspace{1em}
\section{Aim}
\begin{enumerate}
    \item To implement and demonstrate Quantum Fourier Transform (QFT) using the Qiskit framework.
    \item To verify the encoding and decoding of quantum states using QFT and its inverse.
\end{enumerate}

\section{Resources used}
\begin{itemize}
    \item Python
    \item Qiskit 
    \item Jupyter Notebook environment
\end{itemize}

\section{Theoretical Background}
\subsection{Quantum Fourier Transform (QFT)}
The Quantum Fourier Transform is a transformation that maps a quantum state in the computational basis to a superposition of states in the Fourier basis. That is, the Quantum Fourier Transform takes a qubit from the Z-basis to the X-basis. 
\newpage
\subsection{Encoding Steps}
\begin{enumerate}
    \item Prepare the initial quantum state \( |\psi\rangle \) on the first qubit.
    \item Apply a Hadamard gate to the first qubit to create a superposition.
    \item Apply controlled phase shift gates to entangle the qubits and encode the phase information.
    \item Repeat the process for each qubit in the register, applying Hadamard and controlled phase shift gates as necessary.
    \item Finally, apply a series of SWAP gates to reverse the order of the qubits, completing the QFT operation.
\end{enumerate}

\subsection{Decoding Steps}
\begin{enumerate}
    \item Apply a series of SWAP gates to reverse the order of the qubits, preparing for the inverse QFT.
    \item For each qubit, apply controlled phase shift gates in reverse order to disentangle the qubits and decode the phase information.
    \item Apply a Hadamard gate to each qubit to return them to the computational basis.
    \item Measure the qubits to obtain the final state, which should match the original state \( |\psi\rangle \) if the encoding and decoding were successful.
\end{enumerate}

\section{Quantum circuit}
\subsection{Encoding Circuit Setup}
The experiment is conducted following these steps in the Qiskit framework:
\begin{enumerate}
    \item Initialize a quantum circuit with the required number of qubits and classical bits.
    \item Prepare the initial state.
    \item Apply a series of Hadamard and controlled phase shift gates to perform the QFT and obtain the encoded state.
    \item Copy the encoded state to another set of qubits for decoding.
    \item Measure the qubits to obtain the encoded state.
    \item Apply the inverse QFT to the copied qubits to decode the state.
    \item Measure the qubits and record the results.
\end{enumerate}

\begin{figure}[h!]
    \centering
    \includegraphics[width=1.02\textwidth]{../qft_circuit.png}
    \caption{Quantum Fourier Transform Circuit for encoding}
\end{figure}

\newpage
\subsection{Decoding Circuit Setup}
\begin{enumerate}
    \item Apply the inverse QFT to the copied qubits to decode the state.
    \item Measure the qubits and record the results.
\end{enumerate}

\begin{figure}[h!]
    \centering
    \includegraphics[width=1.02\textwidth]{../qft_decode_circuit.png}
    \caption{Quantum Fourier Transform Circuit for decoding}
\end{figure}
\newpage
\section{Measurement Results and Analysis}
The encoding and decoding circuits are executed on the Qiskit Aer Simulator for 2048 shots and also to obtain the bloch sphere representation of the quantum state, and the results obtained are shown in the figures below.

\begin{figure}[h!]
    \centering
    % Remove the \rule line below and uncomment the \includegraphics line to insert your image
    \includegraphics[width=1.0\textwidth]{../input_bloch_vector.png}
    % \rule{0.5\textwidth}{0.3\textwidth} % This creates a placeholder box
    \caption{Bloch sphere representation of the input state}
\end{figure}

\begin{figure}[h!]
    \centering
    % Remove the \rule line below and uncomment the \includegraphics line to insert your image
    \includegraphics[width=1.0\textwidth]{../measurement_encoding.png}
    % \rule{0.5\textwidth}{0.3\textwidth} % This creates a placeholder box
    \caption{Measurement results of the encoded state}
\end{figure}
The measurement results show that the QFT encoding was successful in representing the original state in the combination of superposition states. 

\begin{figure}[h!]
    \centering
    % Remove the \rule line below and uncomment the \includegraphics line to insert your image
    \includegraphics[width=1.0\textwidth]{../output_bloch_vector.png}
    % \rule{0.5\textwidth}{0.3\textwidth} % This creates a placeholder box
    \caption{Bloch sphere representation of the encoded state}
\end{figure}
\newpage

\begin{figure}[h!]
    \centering
    % Remove the \rule line below and uncomment the \includegraphics line to insert your image
    \includegraphics[width=0.5\textwidth]{../measurement_decoding.png}
    % \rule{0.5\textwidth}{0.3\textwidth} % This creates a placeholder box
    \caption{Bloch sphere representation of the decoded state}
\end{figure}
The measurement results show that the inverse QFT decoding was successful in retrieving the original state from the encoded state, confirming the correctness of the QFT implementation.

\section{Conclusion}
The experiment successfully demonstrates the implementation of the Quantum Fourier Transform (QFT) using the Qiskit framework. The encoding and decoding processes were verified through measurement results, confirming that the original quantum state was accurately represented and retrieved.  

\newpage
\section{Appendix}
\subsection{Python Code}
\begin{minted}[fontsize=\small]{python}
import numpy as np
from numpy import pi
from qiskit import QuantumCircuit, transpile
from qiskit.visualization import plot_bloch_multivector, plot_histogram
from qiskit_aer import AerSimulator
import matplotlib.pyplot as plt

def qft_rotations(circuit, n):
    if n == 0:
        return circuit

    n -= 1
    circuit.h(n)
    for qubit in range(n):
        circuit.cp(pi / 2 ** (n - qubit), qubit, n)

    qft_rotations(circuit, n)

def swap_registers(circuit, n):
    for qubit in range(n // 2):
        circuit.swap(qubit, n - qubit - 1)
    return circuit

def qft(circuit, n):
    qft_rotations(circuit, n)
    swap_registers(circuit, n)
    return circuit

n = 4

qc = QuantumCircuit(n)
qft_rotations(qc, n)
fig = qc.draw(output='mpl')
plt.show()

qc = QuantumCircuit(n)
qft(qc, n)
fig = qc.draw(output='mpl')
plt.show()

bits = [0, 1, 1, 0]
qc = QuantumCircuit(n)
for i, bit in enumerate(bits):
    if bit:
        qc.x(i)
fig = qc.draw(output='mpl')
plt.show()

sim = AerSimulator()
qc_init = qc.copy()
qc_init.save_statevector()
state_vector = sim.run(transpile(qc_init, sim)).result().get_statevector()
plot_bloch_multivector(state_vector)
plt.show()

qft(qc, n)
fig = qc.draw(output='mpl')
plt.show()

sim = AerSimulator()
qc_state = qc.copy()
qc_state.save_statevector()
statevector = sim.run(transpile(qc_state, sim)).result().get_statevector()

plot_bloch_multivector(statevector)
plt.show()

qc_qft_measure = qc.copy()
qc_qft_measure.measure_all()
job = sim.run(transpile(qc_qft_measure, sim), shots=2048)
result = job.result().get_counts()
plot_histogram(result, figsize=(15, 6))
plt.show()

inverse_qft = qft(QuantumCircuit(n), n).inverse()
qc.compose(inverse_qft, inplace=True)

qc.decompose()
fig = qc.draw(output='mpl')
plt.show()

qc_final = qc.copy()
qc_final.measure_all()
job1 = sim.run(transpile(qc_final, sim), shots=2048)
result1 = job1.result().get_counts()
plot_histogram(result1, figsize=(7, 6))
plt.show()
\end{minted}
\end{document}
